\documentclass[12pt,a4paper]{article}

\usepackage[utf8]{inputenc}
\usepackage[T1]{fontenc}
\usepackage[polish]{babel}
\usepackage{polski}
\usepackage{geometry}
\geometry{margin=2.5cm}
\usepackage{graphicx}
\usepackage{hyperref}
\usepackage{enumitem}
\usepackage{listings}
\usepackage{booktabs}
\usepackage{float}

% ustawienia wygladu kodu
\lstset{
    basicstyle=\ttfamily\small,
    breaklines=true,
    frame=single,
    captionpos=b,
    extendedchars=true,
    keepspaces=true,
    columns=fixed,
    commentstyle=\normalfont\ttfamily,
    showstringspaces=false,
    literate={ą}{{\k{a}}}1 {ć}{{\'c}}1 {ę}{{\k{e}}}1 {ł}{{\l{}}}1 {ń}{{\'n}}1 {ó}{{\'o}}1 {ś}{{\'s}}1 {ź}{{\'z}}1 {ż}{{\.z}}1 {Ą}{{\k{A}}}1 {Ć}{{\'C}}1 {Ę}{{\k{E}}}1 {Ł}{{\L{}}}1 {Ń}{{\'N}}1 {Ó}{{\'O}}1 {Ś}{{\'S}}1 {Ź}{{\'Z}}1 {Ż}{{\.Z}}1
}

% informacje o autorach
\title{Dokumentacja Implementacyjna Projektu C\\ \large System wyznaczania współrzędnych węzłów grafów planarnych}
\author{Autorzy: Mateusz Bombola, Kajetan Karwacki}
\date{\today}

\begin{document}

\maketitle
\newpage
\tableofcontents
\newpage

\section{Architektura systemu i podział na moduły}
System został zaprojektowany z zachowaniem zasad programowania modułowego. Kod źródłowy w języku C rozdzielono na osobne pliki, oddzielając logikę sterującą od implementacji poszczególnych algorytmów obliczeniowych. Dodatkowo zintegrowano środowisko testowe oparte na skryptach Python.

\subsection{Moduły języka C}
\begin{itemize}
    \item \textbf{\texttt{graph.h}:} Główny plik nagłówkowy spinający projekt. Zawiera definicje struktur danych (\texttt{Vertex}, \texttt{Edge}), stałych matematycznych (np. \texttt{EPSILON}) oraz deklaracje prototypów funkcji obliczeniowych.
    
    \item \textbf{\texttt{main.c}:} Moduł inicjujący i sterujący. Parsuje argumenty linii poleceń , wczytuje i weryfikuje dane topologii, dynamicznie alokuje pamięć, inicjalizuje wstępne współrzędne, a po wykonaniu obliczeń zapisuje wyniki do plików.
    
    \item \textbf{\texttt{smacof.c}:} Niezależny moduł realizujący algorytm SMACOF. Główną logikę umieszczono w funkcji \texttt{calculate\_smacof}, wspieranej przez matematyczną funkcję pomocniczą \texttt{safe\_dist}.
    
    \item \textbf{\texttt{fruchterman.c}:} Moduł zawierający implementację siłowego algorytmu optymalizacyjnego wewnątrz funkcji \texttt{calculate\_fruchterman}.
\end{itemize}

\subsection{Moduły testowe i narzędziowe}
Na potrzeby zapewnienia jakości aplikacji przygotowano zestaw narzędzi:
\begin{itemize}
    \item \textbf{\texttt{generator.py}:} Generator pseudolosowych instancji grafów. Generuje siatki połączeń dla zadanej liczby węzłów i krawędzi (z opcją losowania wag), produkując pliki wejściowe w formacie akceptowalnym przez program.
    \item \textbf{\texttt{wizualizacja.py}:} Analizator graficzny wykorzystujący bibliotekę \texttt{matplotlib}. Wczytuje oryginalną siatkę krawędzi oraz wygenerowane pozycje X i Y, umożliwiając wizualną weryfikację węzłów na płaszczyźnie.
\end{itemize}

\section{Struktury danych}
Podstawowymi strukturami wykorzystywanymi do reprezentacji grafów są:

\begin{lstlisting}[language=C]
typedef struct { 
    int id; 
    double x, y; 
} Vertex;

typedef struct { 
    int u, v; 
    double weight; 
} Edge;
\end{lstlisting}

\textbf{Mechanizm mapowania węzłów (\texttt{main.c}):} Numery węzłów (identyfikatory) podane w pliku wejściowym nie zawsze występują po kolei (np. mogą to być węzły 10 i 100 bez żadnych węzłów pomiędzy). Moduł parsujący buduje strukturę \texttt{node\_map}, przypisując takim węzłom kolejne, ciągłe indeksy ($[0, N-1]$). Pozwala to programowi optymalnie używać pamięci i zapobiega tworzeniu ogromnych, pustych tablic. Dopiero przy zapisie wyników do pliku, program przywraca węzłom ich oryginalne identyfikatory.

\section{Implementacja algorytmów optymalizacyjnych}

\subsection{SMACOF (\texttt{smacof.c})}
Zastosowano algorytm oparty na transformacji Guttmana, składający się z kluczowych faz:
\begin{itemize}
    \item \textbf{Macierz odległości:} Na bazie wczytanych krawędzi algorytm najpierw używa metody \textbf{Floyda-Warshalla} do propagacji najkrótszych ścieżek między wszystkimi niesąsiadującymi węzłami grafu, obsługując izolowane klastry dużą stałą odległością.
    \item \textbf{Inicjalizacja współrzędnych:} Zamiast losować pozycje startowe, program na samym początku układa wszystkie węzły równomiernie na obwodzie koła (korzystając z funkcji \texttt{sin()} i \texttt{cos()}). Zapobiega to nakładaniu się punktów na siebie i pomaga algorytmowi płynniej znaleźć optymalny układ, bez ryzyka utknięcia w złym ułożeniu początkowym.
    \item \textbf{Zabezpieczenie przed dzieleniem przez zero:} Jeśli dwa węzły znajdą się dokładnie w tym samym punkcie (odległość wyniesie 0), wzory algorytmu wygenerowałyby błąd matematyczny psujący wyniki. Aby temu zapobiec, wprowadzono funkcję \texttt{safe\_dist}, która podmienia zerową odległość na minimalną dopuszczalną wartość (\texttt{EPSILON = 1e-9}), zapewniając stabilność całego programu.

\subsection{Fruchterman-Reingold (\texttt{fruchterman.c})}
Zaimplementowano klasyczny model sprężyn i odpychających się ładunków elektrycznych. Obliczenia na każdą iterację \texttt{iterations} dzielą się na:
\begin{itemize}
    \item Faza $O(V^2)$: globalne odpychanie poszczególnych wierzchołków na zasadzie odwrotnej proporcjonalności do ich odległości.
    \item Faza $O(E)$: iteracja wzdłuż bezpośrednich krawędzi (zdefiniowanych w wektorze struktur \texttt{Edge}) i przyciąganie połączonych wierzchołków wzajemnymi wektorami sił z jednoczesnym uwzględnieniem wagi \texttt{(weight)}.
\end{itemize}

\section{Specyfikacja wejścia i wyjścia}
Poniższa specyfikacja definiuje formaty danych wymienianych między modułem obliczeniowym, a zewnętrznymi komponentami systemu.

\subsection{Wejście}
Program parsuje pliki wejściowe liniowo wywołując instrukcję \texttt{fscanf}. Kod zawarty w \texttt{main.c} aplikuje automatyczną naprawę błędnych danych. Dowolne wejściowe wagi krawędzi $\le 0$ są interpretowane i nadpisywane systemową stałą o wartości \texttt{0.1}, co zapobiega propagacji błędu przez ujemne odległości.

\subsection{Format plików tekstowych (\texttt{txt})}
Wynikowa lokalizacja wierzchołków generowana jest do postaci ASCII i powinna być parsowana wyrażeniem: \verb|%d %f %f\n| (Kolejno: Integer ID, Double Oś X, Double Oś Y).

\subsection{Format plików binarnych (\texttt{bin})}
Opcjonalny zrzut strukturalny oparty o procedurę \texttt{fwrite} (\texttt{save\_to\_binary}). Wymaga zastosowania odpowiednich technik deserializacji:
\begin{enumerate}
    \item \textbf{Liczba wierzchołków ($N$):} Pierwsze 4 bajty pliku to zwykła liczba całkowita (typ \texttt{int}). Informuje ona, ile łącznie węzłów zapisano w pliku.
    \item \textbf{Dane pojedynczego wierzchołka:} Reszta pliku natomiast wypełniona jest blokami 20-bajtowymi. Architektura bloku to: $4$ bajty przestrzeni \texttt{id} (Signed Integer 32-bit), $8$ bajtów \texttt{x} (Double 64-bit) oraz $8$ bajtów \texttt{y} (Double 64-bit).
\end{enumerate}

\section{Zarządzanie pamięcią i łączenie}
Kluczowe wektory i macierze przetrzymywane są na stercie. Główna tablica struktur węzłowych \texttt{vertices} jest alokowana z użyciem \texttt{malloc} przez \texttt{main.c} (tylko na podstawie wyliczonej ilości realnie zidentyfikowanych wierzchołków, nie na stałym przedziale rozmiaru grafu). Ponadto tablice dwuwymiarowe w funkcji \texttt{calculate\_smacof} (odległości i wagi) są budowane jako macierze wskaźników dynamicznych.

Wszystkie dynamicznie zaalokowane zasoby są zwalniane za pomocą wywołań funkcji \texttt{free()} po zakończeaniu działania algorytmów lub tuż przed zamknięciem programu, co zapobiega wyciekom pamięci.

\textbf{Komenda kompilacji z wymuszeniem linkowania bibliotek współdzielonych dla trygonometrii:}
\begin{lstlisting}[language=bash]
gcc -o graph_layout main.c smacof.c fruchterman.c -lm
\end{lstlisting}

\end{document}